\documentclass[10pt,a4paper]{article}
\usepackage[T1]{fontenc}
\usepackage[margin=20mm]{geometry}
\usepackage{booktabs,amsmath,listings,xcolor,hyperref}
\hypersetup{colorlinks=true,linkcolor=blue,urlcolor=blue}
\lstset{basicstyle=\ttfamily\small,breaklines=true,columns=fullflexible,
  keepspaces=true,showstringspaces=false,frame=single,rulecolor=\color{gray!35},
  backgroundcolor=\color{gray!4},aboveskip=8pt,belowskip=8pt}
\setlength{\parindent}{0pt}
\setlength{\parskip}{5pt}
\setlength{\emergencystretch}{2em}
\title{SpectralMM: Python/R Usage and Benchmark Reproduction}
\author{}
\date{September 29, 2026}
\begin{document}
% R installation is independent of the standalone CMake build.

\maketitle
\vspace{-1.5em}
This guide covers the shared C++ core and Python/R interfaces. Fitting in Python
or R does not require Julia. Julia is needed to reproduce comparisons using data
shared with the original Julia algorithm. Run all commands from the repository root.

\section*{Automatic Python and Julia installation}
From a local checkout, run \texttt{python -m pip install .}; subsequent Python
sessions use \texttt{import spectralmm}. The wheel contains the shared library.
In Julia, use \texttt{Pkg.develop(path="/path/to/SpectralMM")} and then
\texttt{using SpectralMM}. Call \texttt{SpectralMM.fit(X, y; family=:bernoulli, backend=:cpp)}
for the C++ implementation. Omitting \texttt{backend} selects pure Julia.
Both backends support GLM and non-GLM models through the same matrix API.
The returned model exposes \texttt{model.backend}. On first C++ use, the package builds the core automatically and reuses
its cached library across sessions. All fourteen families are supported.
Existing positional-family Julia methods and \texttt{glm} retain their original
implementation. The manual CMake instructions below are for standalone development.
See \texttt{docs/Automatic\_Installation.md} for details.

\section*{Automatic R package installation}
The repository root is a standard R package. After the updated source is
published, install it with \texttt{remotes::install\_github("Xunjian-Li/SpectralMM")}.
For a local checkout use \texttt{remotes::install\_local("/path/to/SpectralMM")}.
Installation compiles the shared C++ core automatically using R's build system,
RcppEigen and R's configured BLAS. Subsequent sessions only need
\texttt{library(SpectralMM)}. A C++17 toolchain is required for source installation;
R users do not need the standalone CMake commands below. Formula and matrix interfaces are now available in all three languages; see
\texttt{docs/Statistical\_API.md} for the preferred statistical API. See \texttt{docs/R\_Package.md} for the complete workflow.

\section{Files and archived results}
\begin{tabular}{ll}
\toprule
Path & Purpose \\
\midrule
\path{src/} & Original Julia implementation \\
\path{cpp/} & C++ core, C interface and build configuration \\
\path{bindings/python/} & NumPy/SciPy interface \\
\path{R/}, \path{src/} & Standard R package and shared core \\
\path{bindings/julia/} & Native Julia wrapper (Gaussian/Logistic) \\
\path{benchmark/native/} & Data generation, timing and report scripts \\
\path{build/native/} & Local native build, reproducible from source \\
\bottomrule
\end{tabular}

Final results are in \path{benchmark/native/results/family-performance/}.
The \path{report.md} file summarizes findings; CSV files retain per-method results
and individual timings. Configuration and environment records are also retained.
The \path{latex/} directory contains \path{python_tables.tex} and \path{R_tables.tex},
each with four tables covering GLM/non-GLM and dense/sparse cases.
\path{julia_baseline_tables.tex} provides supplementary Julia results.
\path{performance_tables.tex} is a standalone Python/R table document.

Temporary optimization experiments, snapshots and regenerable binary inputs were
removed. Regenerate inputs before rerunning the benchmark (Section~\ref{sec:benchmark}).
The original Julia results in \path{benchmark/scaling_results.csv} and
\path{benchmark/scaling_sparse_results.csv} are also retained.

\section{Dependencies and C++ build}
Building requires a C++17 compiler, CMake and Eigen 3.4 headers. Replace the Eigen
path below with a directory containing \path{Eigen/Core}.

\begin{lstlisting}[language=bash]
export EIGEN3_INCLUDE_DIR=/path/to/eigen-3.4.0
cmake -S cpp -B build/native \
  -DCMAKE_BUILD_TYPE=Release \
  -DEIGEN3_INCLUDE_DIR="$EIGEN3_INCLUDE_DIR" \
  -DSPECTRALMM_USE_BLAS=ON
cmake --build build/native -j2
\end{lstlisting}

The retained binary works on the current machine. Rebuild after changing machines
or source code. This command enables BLAS; the current macOS binary links Apple Accelerate.
Use \texttt{-DSPECTRALMM\_USE\_BLAS=OFF} for an Eigen-only build.
Restart language sessions after replacing a loaded library.

Python requires NumPy; sparse matrices and inference require SciPy. The full
benchmark additionally requires statsmodels and threadpoolctl.

\begin{lstlisting}[language=bash]
python3 -m pip install .
# Optional benchmark dependencies:
python3 -m pip install statsmodels threadpoolctl
\end{lstlisting}

R package installation resolves Rcpp and RcppEigen automatically. Sparse inputs
use Matrix. Standalone benchmarks additionally use jsonlite and MASS.
\begin{lstlisting}[language=R]
install.packages("remotes")
remotes::install_github("Xunjian-Li/SpectralMM")
\end{lstlisting}

The installed Python package loads its bundled library; Julia automatically
builds and reuses a library in its scratch cache. Neither requires a manual path.
Set \texttt{SPECTRALMM\_LIBRARY} to select another standalone build.
The installed R package loads its own registered library automatically.
Filenames are \texttt{libspectralmm.dylib} on macOS, \texttt{libspectralmm.so}
on Linux and \texttt{spectralmm.dll} on Windows.

\section{Python usage}
An intercept is fitted by default. This example supplies six feature columns.
For an existing complete design matrix, set \texttt{fit\_intercept=False} in
Python or \texttt{intercept=FALSE} in R.

\begin{lstlisting}[language=Python]
import numpy as np
from scipy.special import expit
from spectralmm import fit

rng = np.random.default_rng(3124)
Z = rng.normal(size=(200, 6))
X = Z
beta = np.array([0.5, 0.2, -0.3, 0.1, 0.2, -0.1, 0.15])
y = (rng.random(200) < expit(beta[0] + X @ beta[1:])).astype(float)

model = fit(X, y, family="bernoulli", solver="spectral",
            rank=5, floor=1e-6)
print(model.coef)
print(model.info)
probability = model.predict(X)
\end{lstlisting}

\texttt{solver="spectral"} uses spectrally preconditioned CG. The alias \texttt{pcg} selects the same method.
Use \texttt{mm} for spectral MM with a quadratic line search, or
\texttt{cho}/\texttt{cholesky} for a direct weighted least-squares solve.
Cholesky forms a full curvature matrix. Other choices are
\texttt{cg}, \texttt{cgls}, \texttt{crls}, \texttt{lsqr} and \texttt{lsmr}.
SciPy CSC matrices can be passed directly, for example
\texttt{fit(scipy.sparse.csc\_matrix(X), y, ...)}.

Non-GLM models use the same interface, with parameters in \texttt{family\_options}:

\begin{lstlisting}[language=Python]
y_cont = beta[0] + X @ beta[1:] + rng.standard_t(4, size=200)
robust = fit(X, y_cont, family="student_t",
             family_options={"nu": 4, "sigma": 1},
             solver="spectral", floor=1e-6)
fitted = robust.predict(X)

quantile = fit(X, y_cont, family="smooth_quantile",
               family_options={"tau": 0.25, "smoothing": 0.1})
\end{lstlisting}

\section{R usage}
The package installs and registers its native library automatically. Load it once
with \texttt{library(SpectralMM)} before fitting models.

\begin{lstlisting}[language=R]
library(SpectralMM)

set.seed(3124)
Z <- matrix(rnorm(200 * 6), nrow=200)
X <- Z
beta <- c(0.5, 0.2, -0.3, 0.1, 0.2, -0.1, 0.15)
y <- rbinom(200, 1, plogis(as.vector(beta[1] + X %*% beta[-1])))

model <- spectralmm_fit(X, y, family="bernoulli",
                        solver="spectral", rank=5, floor=1e-6)
coef(model)
model$info
probability <- predict(model, X)
\end{lstlisting}

Non-GLM and sparse-input examples:

\begin{lstlisting}[language=R]
y_cont <- as.vector(beta[1] + X %*% beta[-1]) + rt(200, df=4)
robust <- spectralmm_fit(
  X, y_cont, family="student_t",
  family_options=list(nu=4, sigma=1), solver="spectral")
fitted <- predict(robust, X)

library(Matrix)
Xs <- as(Matrix(X, sparse=TRUE), "dgCMatrix")
sparse_model <- spectralmm_fit(Xs, y, family="bernoulli")
\end{lstlisting}

Python and R share one core, but these independent random examples do not generate
identical data. Use the shared binary inputs below when comparing interfaces.

\section{Model names and parameters}
\begin{tabular}{lll}
\toprule
\texttt{family} & Model & Default \texttt{family\_options} \\
\midrule
\texttt{gaussian} & Gaussian identity & None \\
\texttt{bernoulli} & Bernoulli logit & None \\
\texttt{probit} & Bernoulli probit & None \\
\texttt{poisson} & Poisson log & None \\
\texttt{gamma} & Gamma log & None \\
\texttt{negative\_binomial} & Negative-binomial log & \texttt{theta=1} \\
\texttt{smooth\_quantile} & Smooth quantile & \texttt{tau=0.5, smoothing=0.1} \\
\texttt{expectile} & Expectile & \texttt{tau=0.5} \\
\texttt{pseudo\_huber} & Pseudo-Huber & \texttt{delta=1.345} \\
\texttt{student\_t} & Student-$t$ & \texttt{nu=4, sigma=1} \\
\texttt{gaussian\_log} & Gaussian log & None \\
\texttt{gamma\_inverse} & Gamma inverse & None \\
\texttt{tweedie} & Tweedie log & \texttt{power=1.5} \\
\texttt{binomial} & Binomial logit & \texttt{trials=1} \\
\bottomrule
\end{tabular}

\texttt{smoothing} corresponds to Julia's \texttt{epsilon}; \texttt{trials}
can be a scalar or a length-$n$ vector. Binomial predictions return expected counts;
pass \texttt{trials} to \texttt{predict} when new observations have different trial
counts. GLMs return response means; non-GLM models return fitted locations,
quantiles or expectiles. Student-$t$ fitting retains Julia's positive MM surrogate
weights rather than using the exact loss Hessian.

\texttt{floor} is the algorithm's $\rho$, defaulting to $10^{-6}$.
\texttt{ridge} defaults to zero; a nonzero penalty excludes the automatic intercept
unless \texttt{penalize\_intercept} is enabled. Prediction uses the same feature
columns as training. \texttt{beta0} supplies initial coefficients, with the
intercept first when enabled. Benchmarks share initial values across languages.

\section{Optional inference}
By default, inference is attempted for at most 50 coefficients, including the
intercept. Larger models require \texttt{inference=True} in Python or
\texttt{inference=TRUE} in R. False disables inference. The threshold can be
changed with \texttt{inference\_max\_p}. Outputs include standard errors,
t/z statistics, p-values, Wald chi-square statistics and confidence intervals.
GLMs default to model-based covariance; the four residual models default to HC1
sandwich covariance. The optimizer remains matrix-free, but inference forms full
statistical matrices. Numerical and statistical eligibility checks still apply
when inference is explicitly enabled. See \path{docs/README.md} or the complete
\path{SpectralMM_User_Guide.tex} manual for assumptions and accessor examples.

\section{Benchmark reproduction}\label{sec:benchmark}
The archived study covers six GLMs and four non-GLM models with $n=2000$ and
$p=1001$, including the intercept. Four other supported GLM families were not timed.
These commands write to \path{build/family-benchmark/} and preserve archived results.
Ensure \texttt{julia}, \texttt{python3} and \texttt{Rscript} are on PATH.

\begin{lstlisting}[language=bash]
julia --project=. -e 'using Pkg; Pkg.instantiate()'
BENCH_OUT=build/family-benchmark

julia --compiled-modules=existing --project=. \
  benchmark/native/family_performance.jl 1000 "$BENCH_OUT"
\end{lstlisting}

This generates shared data and initial values using the original Julia benchmark
procedure, then warms up and times Julia Spectral-MM. The argument 1000 is the
feature count before adding the intercept; sample size is twice that count.
Run the six native solvers and applicable GLM comparators:

\begin{lstlisting}[language=bash]
python3 benchmark/native/family_performance.py \
  --size 1000 --output "$BENCH_OUT" --language both
\end{lstlisting}

Replace \texttt{both} with \texttt{python} or \texttt{R} to run one interface.
The complete report requires both. GLM allows up to 1000 iterations; nonconverged
Gamma cases can take a long time. Repeat native timings in shuffled order to
reduce execution-order bias:

\begin{lstlisting}[language=bash]
python3 benchmark/native/repeat_family_native.py \
  --output "$BENCH_OUT" --language both
\end{lstlisting}

Generate the Markdown report and LaTeX tables:

\begin{lstlisting}[language=bash]
python3 benchmark/native/report_family_performance.py \
  --output "$BENCH_OUT"
python3 benchmark/native/export_latex.py \
  --results "$BENCH_OUT"

cd "$BENCH_OUT/latex"
pdflatex -interaction=nonstopmode -halt-on-error performance_tables.tex
\end{lstlisting}

Table fragments require \texttt{booktabs}; the standalone document also uses
\texttt{geometry}. Insert \texttt{python\_tables.tex} or \texttt{R\_tables.tex}
into an article. To use the article's final scale, $n=10000$, $p=5001$, replace
1000 with 5000 in both data-generation and Python benchmark commands. That scale
requires more time and memory and differs from the archived run.

\section{Timing and convergence criteria}
Initial runs warm up each method, then report the median of three timings.
Native repeats use five rounds, shuffling six methods each round and batching
short calls. Timings include fitting, wrapper conversion and model construction;
they exclude imports, compilation, input generation and result scoring.
Sparse GLM inputs are densified, with conversion included in GLM timing.
Post-fit inference is disabled for native benchmarks.

Native methods use $\rho=10^{-6}$, no ridge penalty, retained rank $k=5$, initial
Lanczos dimension 12, at most 1000 outer and 200 inner iterations,
\texttt{eta\_max}=0.6, \texttt{correction\_tol}=0.05 and
\texttt{resid\_tol}=0.5, with outer Nesterov acceleration enabled.
Default inner Krylov tolerances are \texttt{rtol}=$10^{-4}$ and
\texttt{atol}=$10^{-8}$. Python/R GLM uses internal tolerance $10^{-12}$;
all results are assessed using the same final gradient criterion.
For gradient $g_0$ at the shared initial coefficients, the strict criterion is
\[
\lVert g\rVert\leq10^{-6}\quad\text{or}\quad
\frac{\lVert g\rVert}{1+\lVert g_0\rVert}\leq10^{-8}.
\]
Default stopping also permits small steps or stagnation. Inspect both
\texttt{converged} and \texttt{gradient\_converged}; solver-declared convergence
alone does not support equal-accuracy speed rankings. A dagger marks declared
convergence without passing the gradient criterion; a double dagger marks binary
separation or perfect prediction.

Archived Logistic/Probit cases exhibit separation, while Gamma GLM reaches the
iteration limit without meeting the criterion. These records are retained;
failures and early stopping are not interpreted as successful-solve speedups.
Main native tables use five-round repeats; GLM and Julia use initial three-run
measurements from different periods. Initial timings showed substantial machine
speed variation, so cross-period differences do not establish language superiority.

\end{document}
